\documentclass[12pt,a4paper]{article}

% ------------------------------------------------------------------
%  Le Codex de l'Architecte — Rapport de projet
%  Compilation : pdflatex RAPPORT.tex  (deux fois pour la table des matières)
% ------------------------------------------------------------------

\usepackage[utf8]{inputenc}
\usepackage[T1]{fontenc}
\usepackage[french]{babel}
\usepackage{lmodern}
\usepackage[a4paper,margin=2.5cm]{geometry}
\usepackage{graphicx}
\usepackage{xcolor}
\usepackage{booktabs}
\usepackage{array}
\usepackage{tabularx}
\usepackage{enumitem}
\usepackage{listings}
\usepackage{titlesec}
\usepackage{fancyhdr}
\usepackage[hidelinks]{hyperref}

% ---- Couleurs du thème -------------------------------------------
\definecolor{cpucyan}{RGB}{0,150,190}
\definecolor{ualorange}{RGB}{200,120,30}
\definecolor{grisdoux}{RGB}{90,100,115}
\definecolor{fondcode}{RGB}{245,247,250}

% ---- En-têtes / pieds de page ------------------------------------
\pagestyle{fancy}
\fancyhf{}
\lhead{\small\itshape Le Codex de l'Architecte}
\rhead{\small\thepage}
\renewcommand{\headrulewidth}{0.4pt}

% ---- Titres de section -------------------------------------------
\titleformat{\section}
  {\normalfont\Large\bfseries\color{cpucyan}}{\thesection}{1em}{}
\titleformat{\subsection}
  {\normalfont\large\bfseries}{\thesubsection}{1em}{}

% ---- Style des blocs de code -------------------------------------
\lstset{
  basicstyle=\ttfamily\small,
  backgroundcolor=\color{fondcode},
  frame=single,
  rulecolor=\color{grisdoux},
  breaklines=true,
  keepspaces=true,
  columns=fullflexible,
  language=,
  showstringspaces=false,
  xleftmargin=1em,
  aboveskip=1em,belowskip=1em
}

% ---- Emplacement pour capture d'écran ----------------------------
% Remplacer le contenu par \includegraphics[width=0.9\linewidth]{fichier}
\newcommand{\capture}[1]{%
  \begin{center}
  \fbox{\parbox[c][3.2cm][c]{0.9\linewidth}{\centering
    \color{grisdoux}\itshape [Emplacement capture d'écran]\\[2mm]#1}}
  \end{center}}

% ==================================================================
\begin{document}

% ------------------------------------------------------------------
%  Page de titre
% ------------------------------------------------------------------
\begin{titlepage}
  \centering
  \vspace*{2cm}
  {\Huge\bfseries Le Codex de l'Architecte\par}
  \vspace{0.6cm}
  {\Large Jeu vidéo éducatif pour l'apprentissage\\[2mm]de la programmation\par}
  \vspace{2cm}
  {\large\itshape Rapport de projet\par}
  \vspace{3cm}
  \begin{tabular}{rl}
    \textbf{Auteur} & : \dots\dots\dots\dots\dots\dots \\[3mm]
    \textbf{Établissement} & : \dots\dots\dots\dots\dots\dots \\[3mm]
    \textbf{Encadrant} & : \dots\dots\dots\dots\dots\dots \\[3mm]
    \textbf{Année académique} & : \dots\dots\dots\dots\dots\dots \\
  \end{tabular}
  \vfill
  {\large Développé avec Unity — Moteur de jeu\par}
\end{titlepage}

% ------------------------------------------------------------------
%  Table des matières
% ------------------------------------------------------------------
\tableofcontents
\thispagestyle{empty}
\clearpage

% ==================================================================
\section{Introduction — Pourquoi ce projet ?}

L'apprentissage de la programmation est l'un des plus grands obstacles rencontrés
par les étudiants en première année d'informatique. Non pas parce que la syntaxe
d'un langage serait insurmontable, mais parce que les \textbf{concepts
fondamentaux restent abstraits} : qu'est-ce qu'une variable ? Où « vit »
réellement une donnée ? Que se passe-t-il, concrètement, lorsqu'on écrit
\texttt{int x = 4;} ou \texttt{z = Int32.Parse(y);} ? Pour beaucoup, le code
reste une suite d'incantations dont le fonctionnement interne demeure invisible.

L'idée de ce projet est née de ce constat. Nous avons voulu créer un outil qui
\textbf{rende visible et tangible} ce qui se passe à l'intérieur de la machine
lorsqu'un programme s'exécute. Plutôt que d'expliquer la mémoire vive avec un
schéma au tableau, pourquoi ne pas laisser l'étudiant \textbf{s'y promener} ?
Plutôt que de décrire le rôle du processeur, pourquoi ne pas l'y faire
\textbf{entrer} ?

C'est ainsi qu'est né \emph{Le Codex de l'Architecte} : un \textbf{jeu vidéo
éducatif en 3D} dans lequel le joueur incarne un agent miniaturisé, parachuté sur
une \textbf{carte mère géante}. Sa mission : exécuter physiquement un programme,
ligne après ligne, en se déplaçant entre les différents composants de
l'ordinateur — la \textbf{RAM} où sont rangées les variables, le \textbf{CPU} qui
lit et calcule, l'\textbf{écran} qui affiche, le \textbf{clavier} qui reçoit la
saisie de l'utilisateur.

L'objectif pédagogique est double :

\begin{itemize}[leftmargin=1.4em]
  \item \textbf{Ancrer les concepts par l'action.} En transportant lui-même une
  valeur du clavier jusqu'à la mémoire, l'étudiant comprend viscéralement la
  différence entre une \emph{variable} (un contenant qui reste en mémoire) et une
  \emph{valeur} (le contenu, qui voyage). Cette distinction, souvent source de
  confusion, devient une évidence physique.
  \item \textbf{Dédramatiser l'erreur et motiver.} Le jeu guide, corrige avec
  bienveillance (l'« humour de l'OS »), récompense la progression par des badges
  et un système de notation, et transforme un cours parfois intimidant en une
  aventure.
\end{itemize}

Le public cible est l'\textbf{étudiant débutant} en informatique, mais l'approche
se veut suffisamment ludique pour intéresser un public plus large et plus jeune.

\capture{L'écran-titre du jeu}

% ==================================================================
\section{Recherche et conception}

Avant d'écrire la première ligne de code, une phase de recherche et de réflexion
a été nécessaire pour poser des bases solides.

\subsection{Étude de l'existant}

Nous avons d'abord étudié les jeux et outils éducatifs qui abordent la
programmation, afin d'en tirer les bonnes idées et d'en éviter les écueils :

\begin{itemize}[leftmargin=1.4em]
  \item \textbf{Human Resource Machine} et \textbf{7 Billion Humans} (Tomorrow
  Corporation) — des puzzles où l'on programme de petits personnages avec des
  instructions élémentaires. Ils démontrent qu'il est possible d'enseigner
  l'algorithmique (variables, boucles, conditions) sans jamais afficher de code
  intimidant, en passant par la manipulation d'objets. Ce fut une inspiration
  majeure pour l'idée de \textbf{manipuler physiquement les données}.
  \item \textbf{Shenzhen I/O} et \textbf{TIS-100} (Zachtronics) — des jeux qui
  plongent le joueur dans l'architecture matérielle (registres, mémoire,
  entrées/sorties). Ils nous ont confortés dans l'idée qu'une \textbf{métaphore
  matérielle} (la carte mère, le CPU, la RAM) est à la fois juste et
  pédagogiquement puissante.
  \item \textbf{CodeCombat}, \textbf{Scratch} — des approches plus scolaires, qui
  montrent l'importance d'une \textbf{progression graduelle} et d'un
  accompagnement constant du débutant.
\end{itemize}

De cette étude, nous avons retenu trois principes directeurs :
\begin{enumerate}[leftmargin=1.6em]
  \item rendre les concepts \textbf{physiques et manipulables} ;
  \item \textbf{cacher la complexité} du langage derrière des actions concrètes ;
  \item \textbf{accompagner} le joueur pas à pas, sans jamais le laisser bloqué.
\end{enumerate}

\subsection{Choix de la plateforme de développement}

Plusieurs moteurs de jeu ont été envisagés :

\begin{center}
\renewcommand{\arraystretch}{1.3}
\begin{tabularx}{\linewidth}{@{}l X l@{}}
\toprule
\textbf{Moteur} & \textbf{Caractéristiques} & \textbf{Décision} \\
\midrule
\textbf{Unity} & Écosystème mûr, langage C\# (cohérent avec le sujet enseigné),
immense documentation, gratuit pour un usage éducatif, multiplateforme &
\textbf{Retenu} \\
Unreal Engine & Rendu photoréaliste, mais surdimensionné et C++ plus lourd &
Écarté \\
Godot & Léger, open source, mais écosystème moins fourni à l'époque du choix &
Écarté \\
\bottomrule
\end{tabularx}
\end{center}

Le choix s'est porté sur \textbf{Unity} (version 6000.4.3f1, dite « Unity 6 »),
avec le pipeline de rendu \textbf{URP} (Universal Render Pipeline) et le nouveau
système d'entrées (\emph{New Input System}). Un argument décisif : Unity utilise
le \textbf{C\#}, c'est-à-dire précisément le langage que le jeu cherche à
enseigner — une cohérence appréciable.

\subsection{Choix des ressources (assets)}

Pour le personnage et ses déplacements, nous avons adopté le pack officiel
\textbf{Starter Assets — ThirdPersonController} de Unity (contrôleur à la
troisième personne, animations, caméra \textbf{Cinemachine}). Cela a permis de
disposer immédiatement d'un personnage jouable et de concentrer l'effort sur le
cœur pédagogique du jeu.

Pour les éléments de décor et les objets manipulés (les boîtes-variables, le
modèle de RAM, le clavier, l'écran plat, le processeur, le portail), des modèles
3D libres ont été importés. \textbf{Un principe fort a toutefois guidé le
développement : tout ce qui pouvait être généré par code l'a été} — l'interface,
le décor de la carte mère, les sons, la musique, la clôture de l'environnement.
Ce choix a rendu le projet \textbf{quasi sans dépendances externes} : il suffit
de cloner le dépôt et d'ouvrir Unity pour que tout fonctionne.

\subsection{Maquettes et validation}

Avant le développement complet, des \textbf{maquettes} (prototypes fonctionnels)
ont été réalisées pour matérialiser l'idée : une première scène où le joueur se
déplace, un premier prototype de « boîte-variable », un embryon de RAM. Ces
maquettes ont été présentées au professeur encadrant afin de \textbf{valider la
direction} avant d'investir des mois de développement. Cette validation obtenue,
le développement à proprement parler a pu commencer.

\capture{Une des premières maquettes}

% ==================================================================
\section{Développement du jeu}

Cette partie constitue le cœur du projet. Elle s'est étalée sur \textbf{plusieurs
mois} et représente la charge de travail la plus importante. Nous la détaillerons
selon quatre axes : l'architecture technique, le concept pédagogique central, la
structure du jeu, et enfin le contenu.

\subsection{Architecture technique générale}

Le projet est organisé autour d'un \textbf{script central, \texttt{GameState}},
implémenté en \emph{singleton} persistant (patron \emph{singleton} +
\texttt{DontDestroyOnLoad}). Ce composant survit aux changements de scène et
centralise \textbf{tout l'état du jeu} : ce que le joueur porte, le contenu de la
mémoire, la progression dans le programme, le score, la sauvegarde. Toutes les
autres briques du jeu consultent et modifient cet état unique, ce qui évite les
incohérences.

Autour de ce noyau gravitent des \textbf{composants autonomes}, chacun
responsable d'un système précis et créé automatiquement au lancement :
interface (HUD, marqueur d'objectif, journal, notifications, menu pause,
minimap), audio (voix radio, humour de l'OS, musique), décor (clôture, composants
de carte mère, easter egg) et aide (le drone d'assistance).

Chaque système suit le même patron : une méthode statique \texttt{Ensure()} qui
le crée s'il n'existe pas, et un \texttt{Awake()} qui garantit l'unicité. Cette
régularité rend le code \textbf{prévisible et facile à étendre} — un point
important pour un projet destiné à être repris et poursuivi.

Le jeu est découpé en \textbf{cinq scènes} reliées entre elles :

\begin{center}
\renewcommand{\arraystretch}{1.3}
\begin{tabularx}{\linewidth}{@{}l X@{}}
\toprule
\textbf{Scène} & \textbf{Rôle} \\
\midrule
\texttt{MainScene} & Le monde ouvert (la carte mère) où le joueur se déplace \\
\texttt{RAM} & La mémoire vive : déclaration, lecture et écriture des variables \\
\texttt{CPU} & L'\textbf{unité de contrôle} : lit le programme, révèle les instructions \\
\texttt{Calculateur} & L'\textbf{unité arithmétique et logique} : conversions, calculs, tests \\
\texttt{Clavier} & Le terminal de saisie (\texttt{Console.ReadLine}) \\
\bottomrule
\end{tabularx}
\end{center}

\capture{Vue d'ensemble de la carte mère (MainScene)}

\subsection{Le concept pédagogique central : variable \texorpdfstring{$\neq$}{different de} valeur}

Toute la conception du gameplay repose sur une \textbf{distinction fondamentale},
exigée par la rigueur pédagogique :

\begin{center}
\renewcommand{\arraystretch}{1.3}
\begin{tabularx}{\linewidth}{@{}l X X@{}}
\toprule
 & \textbf{Variable (une boîte)} & \textbf{Valeur (un contenu)} \\
\midrule
Où vit-elle ? & Dans la RAM, qu'elle ne quitte jamais & Elle voyage sur la tête du joueur \\
Représentation & Une boîte en carton (nom, type, valeur) & Juste le texte, coloré selon le type \\
Naissance & Par une déclaration (type, nom) & Lecture, saisie clavier, ou calcul \\
Lire & Cliquer la boîte → copie de sa valeur & — \\
Écrire (affectation) & Arriver avec une valeur, cliquer la variable & — \\
\bottomrule
\end{tabularx}
\end{center}

Ce modèle rend concrets des mécanismes qui restent d'ordinaire abstraits : une
\textbf{déclaration} réserve une case mémoire ; \textbf{lire} une variable ne la
détruit pas (on n'emporte qu'une copie) ; une \textbf{affectation} remplace le
contenu ; et le \textbf{CPU ne manipule que des valeurs}, jamais des boîtes — car
un processeur calcule sur des données, pas sur des emplacements mémoire.

Chaque \textbf{type} possède sa couleur (\texttt{int} en rouge, \texttt{float} en
magenta, \texttt{string} en bleu, \texttt{bool} en noir, \texttt{char} en bleu
clair), et cette couleur suit la valeur partout où elle va.

\capture{Le joueur portant une valeur ; une variable dans la RAM}

\subsection{Le CPU en deux unités}

Fidèle à l'architecture réelle d'un processeur, le CPU est scindé en \textbf{deux
unités visuellement distinctes} :

\begin{itemize}[leftmargin=1.4em]
  \item l'\textbf{\textcolor{cpucyan}{unité de contrôle}} — c'est là que le joueur
  \textbf{lit le programme} ; chaque instruction lui est révélée une à une,
  l'obligeant à revenir consulter le CPU avant d'agir (comme un processeur qui
  charge l'instruction suivante) ;
  \item l'\textbf{\textcolor{ualorange}{unité arithmétique et logique (UAL)}} —
  c'est là que s'effectuent les \textbf{conversions, additions et tests}. Fait
  notable : l'UAL \textbf{ignore la provenance et la destination} des données ;
  elle affiche seulement \texttt{... + ...}, jamais \texttt{somme = ...}, car un
  additionneur ne fait qu'additionner — c'est le programme qui sait où ira le
  résultat.
\end{itemize}

Cette séparation, matérialisée par des couleurs, des bandeaux et des filigranes
opposés, ancre une notion d'architecture des ordinateurs souvent négligée.

\capture{L'unité de contrôle et l'unité arithmétique}

\subsection{Le programme du chapitre 1 (les six lignes de base)}

Le premier chapitre fait exécuter au joueur un programme complet de six lignes,
qui couvre les fondamentaux :

\begin{center}
\renewcommand{\arraystretch}{1.3}
\begin{tabularx}{\linewidth}{@{}l X@{}}
\toprule
\textbf{Ligne} & \textbf{Concept enseigné} \\
\midrule
\texttt{int x = 4;} & La \textbf{déclaration} d'une variable \\
\texttt{Console.WriteLine(x);} & L'\textbf{affichage} d'une valeur \\
\texttt{string y = Console.ReadLine();} & La \textbf{saisie} utilisateur au clavier \\
\texttt{int z = Int32.Parse(y);} & La \textbf{conversion} de type (texte $\rightarrow$ entier) \\
\texttt{int somme = x + z;} & Le \textbf{calcul} arithmétique \\
\texttt{Console.WriteLine(somme);} & L'affichage du résultat \\
\bottomrule
\end{tabularx}
\end{center}

Chaque ligne se décompose en \textbf{étapes} que le joueur franchit physiquement.
Par exemple, la ligne 4 (\texttt{int z = Int32.Parse(y);}) se joue ainsi :
(1)~déclarer la variable \texttt{z} dans la RAM ; (2)~aller lire la valeur de
\texttt{y} et l'apporter à l'unité arithmétique, qui la convertit et rend un
entier ; (3)~revenir ranger cet entier dans \texttt{z}. Le joueur \emph{vit} la
conversion au lieu de la lire.

Un détail pédagogique important : à la ligne 3, la valeur récupérée au clavier
arrive \textbf{sans nom} — c'est seulement en la rangeant dans \texttt{y} qu'elle
acquiert son identité, illustrant que \texttt{y = Console.ReadLine()} est une
\textbf{affectation}.

\capture{Le programme affiché sur le CPU}

\subsection{Le chapitre 2 : la condition \texttt{if} — les deux portes}

Le second chapitre introduit la \textbf{structure conditionnelle} avec une mise
en scène marquante. La ligne à exécuter est :

\begin{lstlisting}
if (somme > 50) { Console.WriteLine("grand"); }
else            { Console.WriteLine("petit"); }
\end{lstlisting}

Le déroulé pédagogique :
\begin{enumerate}[leftmargin=1.6em]
  \item Le joueur apporte la valeur de \texttt{somme} à l'unité arithmétique, qui
  \textbf{évalue le test} sous ses yeux (\texttt{67 > 50} $\rightarrow$ VRAI) et
  lui rend un \textbf{booléen} — matérialisant que le résultat d'une condition est
  une valeur de type \texttt{bool}.
  \item De retour dans le monde, \textbf{deux portes} se dressent devant l'écran :
  l'une pour la branche \texttt{VRAI} $\rightarrow$ « grand », l'autre pour
  \texttt{FAUX} $\rightarrow$ « petit ».
  \item \textbf{Seule la porte correspondant au booléen s'ouvre.} L'autre est
  murée : s'en approcher rappelle qu'une branche dont le test échoue ne s'exécute
  jamais.
  \item Traverser la bonne porte \textbf{exécute la branche} : l'écran affiche le
  message.
\end{enumerate}

Comme la valeur de \texttt{somme} dépend du nombre aléatoire saisi au chapitre 1,
\textbf{chaque partie emprunte une branche différente} — la meilleure
démonstration possible de ce qu'est une condition.

\capture{Les deux portes de la condition}

\subsection{Les systèmes annexes}

Autour de ce cœur pédagogique, de nombreux systèmes enrichissent l'expérience :

\begin{itemize}[leftmargin=1.4em]
  \item \textbf{La voix radio.} Un « centre de contrôle » guide le joueur à chaque
  instruction, avec une \textbf{voix neuronale de synthèse} (générée hors ligne),
  bien plus naturelle qu'une voix robotique. Les consignes ne se déclenchent qu'à
  la sortie du CPU, une fois l'instruction lue.
  \item \textbf{Le guidage et le verrouillage.} Le joueur ne peut pas « sauter »
  une étape : tant qu'il n'a pas lu la prochaine instruction au CPU, les stations
  refusent d'interagir. S'il se trompe de station, un message le réoriente. Le
  HUD, un marqueur d'objectif 3D et une minimap indiquent où aller.
  \item \textbf{Le drone d'aide.} Un drone flottant, activable à la demande,
  explique le concept de la ligne en cours et rappelle les règles.
  \item \textbf{La progression.} Un système de \textbf{badges} (un par notion), de
  \textbf{notation} par ligne et de \textbf{skins} débloquables récompense
  l'apprentissage. Un \textbf{rapport élève} exportable permet au professeur de
  suivre les progrès.
  \item \textbf{L'ambiance.} Écran-titre avec caméra orbitale, cinématique
  d'introduction, musique composée par code, décor procédural de carte mère,
  clôture de l'environnement, et un \textbf{easter egg pédagogique} : un insecte
  caché rappelant le premier « bug » de l'histoire de l'informatique (Grace
  Hopper, 1947).
  \item \textbf{L'accessibilité.} Un « mode Zen » retire la pression du
  chronomètre, les messages d'erreur sont bienveillants, et l'ensemble est en
  français.
\end{itemize}

\capture{Le HUD, la minimap, une notification, le drone d'aide}

\subsection{Sauvegarde et robustesse}

La progression du joueur (ligne courante, contenu de la mémoire, score, badges,
réglages) est \textbf{sauvegardée automatiquement}. Le joueur peut quitter et
\textbf{reprendre exactement où il s'était arrêté}, sans revoir la cinématique.
Une grande attention a été portée à la robustesse : gestion des changements de
scène, périodes de « grâce » pour éviter des messages intempestifs, compatibilité
des anciennes sauvegardes lorsque le programme s'allonge.

% ==================================================================
\section{Difficultés rencontrées}

Le développement, étalé sur plusieurs mois, n'a pas été sans obstacles :

\begin{itemize}[leftmargin=1.4em]
  \item \textbf{La traduction d'un concept abstrait en gameplay.} La distinction
  variable/valeur, évidente sur le papier, a demandé de \textbf{nombreuses
  itérations} avant de trouver sa forme jouable. La solution — les variables
  restent en RAM, seules les valeurs voyagent — n'a émergé qu'après plusieurs
  refontes.
  \item \textbf{La gestion de l'état entre cinq scènes.} Faire communiquer
  proprement le monde, la RAM, le CPU, l'UAL et le clavier, tout en gardant un
  état cohérent, a nécessité une architecture centralisée et de nombreux
  ajustements (synchronisation, objets détruits au changement de scène).
  \item \textbf{Les enchaînements audio et cinématiques.} Faire coïncider
  l'écran-titre, la cinématique, la musique et la voix radio sans qu'ils se
  chevauchent a demandé la mise en place d'une file d'attente vocale et d'un
  séquencement précis.
  \item \textbf{L'adaptation à l'affichage.} Le passage en plein écran révélait
  des problèmes de cadrage, corrigés par une adaptation dynamique de la caméra.
  \item \textbf{Des problèmes techniques ponctuels} : un fichier source Blender
  bloquant l'import, des chevauchements de texte dans l'interface, des messages
  mal synchronisés — autant de petits obstacles résolus un à un.
\end{itemize}

Chacune de ces difficultés a été l'occasion d'\textbf{apprendre} : sur Unity, sur
l'architecture logicielle, et sur la conception pédagogique.

% ==================================================================
\section{Points d'amélioration}

Le jeu, pleinement fonctionnel pour son objectif (le chapitre 1 et le début du
chapitre 2), offre de nombreuses perspectives d'évolution :

\begin{itemize}[leftmargin=1.4em]
  \item \textbf{Le design visuel.} L'esthétique actuelle, largement générée par
  code, est fonctionnelle mais perfectible. Un travail artistique (textures,
  éclairage, effets) rendrait l'univers plus immersif.
  \item \textbf{Le chapitre 2 et au-delà.} La condition \texttt{if} est en place ;
  il reste à développer pleinement les \textbf{boucles} (\texttt{for},
  \texttt{while}), dont la logique est déjà présente en réserve dans le code.
  Au-delà, on pourrait aborder les tableaux, les fonctions, les objets.
  \item \textbf{Enrichissement pédagogique.} Niveaux « bac à sable », défis
  chronométrés, mode « débogueur » pas à pas.
  \item \textbf{Diffusion.} Un export \textbf{WebGL} permettrait de jouer dans un
  navigateur, sans installation — idéal en salle de classe.
  \item \textbf{Accessibilité et internationalisation.} Sous-titres, mode
  daltonien, traduction en plusieurs langues.
\end{itemize}

Ces pistes sont documentées afin qu'un futur développeur puisse \textbf{reprendre
et poursuivre} le projet aisément.

% ==================================================================
\section{Conclusion}

\emph{Le Codex de l'Architecte} est né d'une conviction : \textbf{on apprend mieux
en faisant}. En transformant les concepts abstraits de la programmation —
variables, mémoire, conversion, condition — en actions physiques au sein d'une
carte mère géante, le jeu offre à l'étudiant débutant une compréhension
\textbf{intuitive et durable} de ce qui se passe réellement lorsqu'un programme
s'exécute.

Le projet a représenté un travail conséquent, étalé sur plusieurs mois, mêlant
\textbf{conception pédagogique}, \textbf{développement logiciel} et
\textbf{création d'expérience}. Au-delà du résultat — un jeu jouable, guidé,
récompensant —, il aura été une formidable occasion d'apprentissage : de Unity et
du C\#, de l'architecture d'une application complète, et de l'art délicat de
rendre l'abstrait tangible.

Le socle est solide, l'architecture pensée pour l'extension, et les fondations du
chapitre 2 posées. \emph{Le Codex de l'Architecte} ne demande qu'à grandir —
chapitre après chapitre — pour accompagner toujours plus loin celles et ceux qui
font leurs premiers pas dans le monde du code.

\capture{L'écran de fin / le rating}

% ==================================================================
\section{Références}

\subsection*{Documentation technique}
\begin{itemize}[leftmargin=1.4em]
  \item \textbf{Unity — Documentation officielle}, Unity Technologies.
  \url{https://docs.unity3d.com}
  \item \textbf{Unity — Universal Render Pipeline (URP)}.
  \url{https://docs.unity3d.com/Packages/com.unity.render-pipelines.universal@latest}
  \item \textbf{Unity — Input System}.
  \url{https://docs.unity3d.com/Packages/com.unity.inputsystem@latest}
  \item \textbf{Unity — Starter Assets: ThirdPerson}, Unity Asset Store.
  \item \textbf{Unity — Cinemachine}, documentation.
  \item \textbf{Microsoft — Documentation C\#} (types, \texttt{Console},
  \texttt{Int32.Parse}). \url{https://learn.microsoft.com/dotnet/csharp}
  \item \textbf{edge-tts} — synthèse vocale neuronale.
  \url{https://pypi.org/project/edge-tts}
\end{itemize}

\subsection*{Jeux et outils ayant inspiré le projet}
\begin{itemize}[leftmargin=1.4em]
  \item \textbf{Human Resource Machine} / \textbf{7 Billion Humans} — Tomorrow
  Corporation.
  \item \textbf{Shenzhen I/O} / \textbf{TIS-100} — Zachtronics.
  \item \textbf{CodeCombat} — CodeCombat Inc.
  \item \textbf{Scratch} — MIT Media Lab.
\end{itemize}

\subsection*{Culture informatique}
\begin{itemize}[leftmargin=1.4em]
  \item L'anecdote du premier « bug » (papillon dans le Harvard Mark II, équipe de
  \textbf{Grace Hopper}, 1947), reprise comme easter egg pédagogique.
\end{itemize}

\end{document}
